% Concise project summary snippet for thesis report
% You can % Concise project summary snippet for thesis report
% You can % Concise project summary snippet for thesis report
% You can % Concise project summary snippet for thesis report
% You can \input{report_code_summary.tex} in your main .tex document.

\section{Conceptual Heat Pump Model With Turbine Expansion}

\textbf{Objective.} The objective of this code is to quantify \textbf{how much system COP can increase} when replacing throttling losses with a \textbf{turbine expander} and mechanical work recovery, using a global optimization of cycle state variables.

\subsection{Modelling Approach}

The implementation solves a steady-state vapor-compression heat pump cycle with turbine expansion, pinch-point constraints, and performance post-processing. Thermophysical properties are evaluated with \textbf{CoolProp/REFPROP} through an AbstractState-based wrapper. The optimization is performed with \textbf{Differential Evolution} and candidate cycles are parameterized by \(PR,\,T_{ref,1},\,s_{ref,1},\,T_{ref,3}\) (or \(T_{ref,1},\,s_{ref,1},\,T_{ref,3}\) when \(PR\) is fixed).

\textbf{Key assumptions}
\begin{itemize}
\item Steady-state, 1D component model.
\item Fixed component efficiencies: $\eta_{\mathrm{compr}}$, $\eta_{\mathrm{turb}}$, and shaft efficiency $\eta_{\mathrm{shaft}}$.
\item Pinch constraints are enforced as \textbf{true minimum approach temperatures} along condenser and evaporator heat-transfer paths (not only at endpoints).
\item Refrigerant mass flow rate is determined from the required sink-side heat duty $\dot Q_{out,req}$ and candidate specific condenser heat rejection.
\item Subcritical and supercritical operating regions both handled, with a near-critical buffer to avoid numerically unreliable states.
\item Infeasible candidates receive a small positive penalty score, while feasible candidates are ranked by minimizing $-COP_{\mathrm{turb}}$.
\end{itemize}

\textbf{Block diagram (code workflow)}

\textit{(Requires in preamble: \texttt{\textbackslash usepackage\{tikz\}} and
\texttt{\textbackslash usetikzlibrary\{arrows.meta,positioning,shapes.geometric\}}.)}

\begin{figure}[h!]
\centering
\begin{tikzpicture}[
	node distance=1.25cm and 1.8cm,
	>=Latex,
	block/.style={draw, rounded corners, align=left, text width=0.28\textwidth, inner sep=6pt, font=\footnotesize},
	decision/.style={draw, diamond, aspect=2.0, align=center, text width=0.18\textwidth, inner sep=2pt, font=\footnotesize},
	line/.style={-Latex, thick}
]

\node[block] (in) {
\textbf{Inputs / configuration}\\
$\{T_{h,in},\,T_{c,in},\,\dot m_h,\,\dot m_c,\,c_{p,h},\,c_{p,c},\,\eta_{\mathrm{compr}},\,\eta_{\mathrm{turb}},\,\eta_{\mathrm{shaft}},\\
\Delta T_{\mathrm{pp},1},\,\Delta T_{\mathrm{pp},3},\,\Delta T_{\mathrm{pp},4},\,\Delta T_{\mathrm{sh}},\,\text{refrigerant}\}$
};

\node[decision, right=of in] (mode) {\textbf{Fixed $PR$?}};

\node[block, above right=0.35cm and 1.7cm of mode] (eval) {
\textbf{Cycle evaluation for candidate variables}\\
Solve stations $1\rightarrow2\rightarrow3\rightarrow4$:\\
$\{p_i,T_i,h_i,s_i,Q_i\}$ from CoolProp/REFPROP.
};

\node[block, below right=0.35cm and 1.7cm of mode] (obj) {
\textbf{Objective function $f(\mathbf{x})$}\\
$\mathbf{x}=[PR,\,T_{ref,1},\,s_{ref,1},\,T_{ref,3}]$ (or fixed-$PR$ subset),\\
$f(\mathbf{x})=-COP_{\mathrm{turb}}$ with infeasibility penalty.
};

\node[block, right=of eval] (checks) {
\textbf{Feasibility checks}\\
Pinch constraints: true condenser/evaporator $\Delta T_{\min}$ checks,\\
physical checks (quality/state validity),\\
near-critical buffer for robust optimization.
};

\node[block, right=of obj] (opt) {
\textbf{Global optimization (DE)}\\
Bounded search over candidate vector $\mathbf{x}$;\\
select $\mathbf{x}^*=\arg\max(COP_{\mathrm{turb}})$ among feasible solutions.
};

\node[block, right=of checks] (perf) {
\textbf{Performance post-processing}\\
Compute $\dot W_{\mathrm{compr}},\,\dot W_{\mathrm{turb}},\,\dot Q_{out},\\
COP_{\mathrm{turb}},\,COP_{\mathrm{is}},\,COP_{\mathrm{isenth}}$.
};

\node[block, below=1.1cm of perf] (out) {
\textbf{Outputs}\\
Best cycle state set, selected decision variables, COP metrics, and TS/PH plotting data.
};

\draw[line] (in) -- (mode);
\draw[line] (mode) -- node[above, font=\scriptsize] {yes} (eval);
\draw[line] (mode) -- node[below, font=\scriptsize] {no} (obj);
\draw[line] (eval) -- (checks);
\draw[line] (checks) -- (perf);
\draw[line] (obj) -- (opt);
\draw[line] (opt) |- (perf);
\draw[line] (perf) -- (out);

\end{tikzpicture}
\caption{\textbf{Detailed workflow of the turbine-assisted heat pump model and global multi-variable optimization.}}
\end{figure}

\subsection{Implemented Optimization Method (Current Version)}

The current solver uses \textbf{Differential Evolution} (global, population-based search) instead of local or pressure-ratio-only scanning. The candidate bounds are defined as:
\begin{itemize}
\item $PR \in [1.01,\,30]$ (when not fixed),
\item $T_{ref,1} \in [T_{c,in}-30,\,T_{c,in}-\Delta T_{pp,min,evap}]$,
\item $s_{ref,1}$ around saturated-vapor entropy at the cold-side inlet condition,
\item $T_{ref,3} \in [T_{h,in}+\Delta T_{pp,min,cond},\,T_{h,in}+150]$.
\end{itemize}

For each candidate, the model reconstructs all thermodynamic states, computes the required refrigerant mass flow from $\dot Q_{out,req}$, and rejects nonphysical or pinch-violating cycles. The best feasible solution maximizes $COP_{\mathrm{turb}}$ (equivalently minimizes $-COP_{\mathrm{turb}}$).

\subsection{Verification}

The implementation is verified against the model of \textbf{Martin T. White (PocketTHERM / SoftwareX)}. The reported agreement is within \textbf{0.01\% maximum relative error}. This supports correctness of the implemented cycle equations and property-coupled solution workflow.

\subsection{Cycle Input Case Used in Report (R1234ze(E))}

\begin{table}[h!]
\centering
\caption{\textbf{Cycle input parameters for the reported R1234ze(E) case}.}
\begin{tabular}{lll}
\hline
\textbf{Parameter} & \textbf{Value} & \textbf{Unit} \\
\hline
Refrigerant & R1234ze(E) & -- \\
$T_{h,in}$ & 353.15 & K \\
$T_{c,in}$ & 287.15 & K \\
$\dot m_h$ & 42 & kg/s \\
$c_{p,h}$ & 1885 & J/(kg\,K) \\
$\dot m_c$ & 40 & kg/s \\
$c_{p,c}$ & 1006 & J/(kg\,K) \\
$\eta_{\mathrm{compr}}$ & 0.78 & -- \\
$\eta_{\mathrm{turb}}$ & 0.87 & -- \\
$\eta_{\mathrm{shaft}}$ & 0.98 & -- \\
$\Delta T_{\mathrm{pp},1}$ & 10 & K \\
$\Delta T_{\mathrm{pp},3}$ & 10 & K \\
$\Delta T_{\mathrm{pp},4}$ & 10 & K \\
$\Delta T_{\mathrm{sh}}$ & 5 & K \\
\hline
\end{tabular}
\end{table}

\subsection{Cycle Input Case Used for CO2 Analyses}

\begin{table}[h!]
\centering
\caption{\textbf{Cycle input parameters used for CO2 analyses}.}
\begin{tabular}{lll}
\hline
\textbf{Parameter} & \textbf{Value} & \textbf{Unit} \\
\hline
Refrigerant & CO2 & -- \\
$T_{h,in}$ & 276 & K \\
$T_{c,in}$ & 240 & K \\
$\dot m_h$ & 42 & kg/s \\
$c_{p,h}$ & 1885 & J/(kg\,K) \\
$\dot m_c$ & 80 & kg/s \\
$c_{p,c}$ & 1006 & J/(kg\,K) \\
$\eta_{\mathrm{compr}}$ & 0.78 & -- \\
$\eta_{\mathrm{turb}}$ & 0.87 & -- \\
$\eta_{\mathrm{shaft}}$ & 0.98 & -- \\
$\Delta T_{\mathrm{pp},1}$ & 10 & K \\
$\Delta T_{\mathrm{pp},3}$ & 10 & K \\
$\Delta T_{\mathrm{pp},4}$ & 3 & K \\
$\Delta T_{\mathrm{sh}}$ & 5 & K \\
\hline
\end{tabular}
\end{table}

\subsection{Cycle COP Results (R1234ze(E) Case)}

\begin{table}[h!]
\centering
\caption{\textbf{Computed performance results for the R1234ze(E) case}.}
\begin{tabular}{lll}
\hline
\textbf{Metric} & \textbf{Value} & \textbf{Unit} \\
\hline
Optimized pressure ratio $PR$ & \textbf{12.684} & -- \\
$COP_{\mathrm{turb}}$ & \textbf{2.690} & -- \\
$COP_{\mathrm{is}}$ & 2.875 & -- \\
$COP_{\mathrm{isenth}}$ & 1.879 & -- \\
$\dot Q_{out}$ & 318492 & W \\
$\dot W_{compr}$ & 169529 & W \\
$\dot W_{turb}$ & 52171 & W \\
\hline
\end{tabular}
\end{table}

\textbf{Interpretation.} The model predicts a substantial COP uplift potential when turbine work recovery is included, while satisfying the imposed pinch constraints.
 in your main .tex document.

\section{Conceptual Heat Pump Model With Turbine Expansion}

\textbf{Objective.} The objective of this code is to quantify \textbf{how much system COP can increase} when replacing throttling losses with a \textbf{turbine expander} and mechanical work recovery, using a global optimization of cycle state variables.

\subsection{Modelling Approach}

The implementation solves a steady-state vapor-compression heat pump cycle with turbine expansion, pinch-point constraints, and performance post-processing. Thermophysical properties are evaluated with \textbf{CoolProp/REFPROP} through an AbstractState-based wrapper. The optimization is performed with \textbf{Differential Evolution} and candidate cycles are parameterized by \(PR,\,T_{ref,1},\,s_{ref,1},\,T_{ref,3}\) (or \(T_{ref,1},\,s_{ref,1},\,T_{ref,3}\) when \(PR\) is fixed).

\textbf{Key assumptions}
\begin{itemize}
\item Steady-state, 1D component model.
\item Fixed component efficiencies: $\eta_{\mathrm{compr}}$, $\eta_{\mathrm{turb}}$, and shaft efficiency $\eta_{\mathrm{shaft}}$.
\item Pinch constraints are enforced as \textbf{true minimum approach temperatures} along condenser and evaporator heat-transfer paths (not only at endpoints).
\item Refrigerant mass flow rate is determined from the required sink-side heat duty $\dot Q_{out,req}$ and candidate specific condenser heat rejection.
\item Subcritical and supercritical operating regions both handled, with a near-critical buffer to avoid numerically unreliable states.
\item Infeasible candidates receive a small positive penalty score, while feasible candidates are ranked by minimizing $-COP_{\mathrm{turb}}$.
\end{itemize}

\textbf{Block diagram (code workflow)}

\textit{(Requires in preamble: \texttt{\textbackslash usepackage\{tikz\}} and
\texttt{\textbackslash usetikzlibrary\{arrows.meta,positioning,shapes.geometric\}}.)}

\begin{figure}[h!]
\centering
\begin{tikzpicture}[
	node distance=1.25cm and 1.8cm,
	>=Latex,
	block/.style={draw, rounded corners, align=left, text width=0.28\textwidth, inner sep=6pt, font=\footnotesize},
	decision/.style={draw, diamond, aspect=2.0, align=center, text width=0.18\textwidth, inner sep=2pt, font=\footnotesize},
	line/.style={-Latex, thick}
]

\node[block] (in) {
\textbf{Inputs / configuration}\\
$\{T_{h,in},\,T_{c,in},\,\dot m_h,\,\dot m_c,\,c_{p,h},\,c_{p,c},\,\eta_{\mathrm{compr}},\,\eta_{\mathrm{turb}},\,\eta_{\mathrm{shaft}},\\
\Delta T_{\mathrm{pp},1},\,\Delta T_{\mathrm{pp},3},\,\Delta T_{\mathrm{pp},4},\,\Delta T_{\mathrm{sh}},\,\text{refrigerant}\}$
};

\node[decision, right=of in] (mode) {\textbf{Fixed $PR$?}};

\node[block, above right=0.35cm and 1.7cm of mode] (eval) {
\textbf{Cycle evaluation for candidate variables}\\
Solve stations $1\rightarrow2\rightarrow3\rightarrow4$:\\
$\{p_i,T_i,h_i,s_i,Q_i\}$ from CoolProp/REFPROP.
};

\node[block, below right=0.35cm and 1.7cm of mode] (obj) {
\textbf{Objective function $f(\mathbf{x})$}\\
$\mathbf{x}=[PR,\,T_{ref,1},\,s_{ref,1},\,T_{ref,3}]$ (or fixed-$PR$ subset),\\
$f(\mathbf{x})=-COP_{\mathrm{turb}}$ with infeasibility penalty.
};

\node[block, right=of eval] (checks) {
\textbf{Feasibility checks}\\
Pinch constraints: true condenser/evaporator $\Delta T_{\min}$ checks,\\
physical checks (quality/state validity),\\
near-critical buffer for robust optimization.
};

\node[block, right=of obj] (opt) {
\textbf{Global optimization (DE)}\\
Bounded search over candidate vector $\mathbf{x}$;\\
select $\mathbf{x}^*=\arg\max(COP_{\mathrm{turb}})$ among feasible solutions.
};

\node[block, right=of checks] (perf) {
\textbf{Performance post-processing}\\
Compute $\dot W_{\mathrm{compr}},\,\dot W_{\mathrm{turb}},\,\dot Q_{out},\\
COP_{\mathrm{turb}},\,COP_{\mathrm{is}},\,COP_{\mathrm{isenth}}$.
};

\node[block, below=1.1cm of perf] (out) {
\textbf{Outputs}\\
Best cycle state set, selected decision variables, COP metrics, and TS/PH plotting data.
};

\draw[line] (in) -- (mode);
\draw[line] (mode) -- node[above, font=\scriptsize] {yes} (eval);
\draw[line] (mode) -- node[below, font=\scriptsize] {no} (obj);
\draw[line] (eval) -- (checks);
\draw[line] (checks) -- (perf);
\draw[line] (obj) -- (opt);
\draw[line] (opt) |- (perf);
\draw[line] (perf) -- (out);

\end{tikzpicture}
\caption{\textbf{Detailed workflow of the turbine-assisted heat pump model and global multi-variable optimization.}}
\end{figure}

\subsection{Implemented Optimization Method (Current Version)}

The current solver uses \textbf{Differential Evolution} (global, population-based search) instead of local or pressure-ratio-only scanning. The candidate bounds are defined as:
\begin{itemize}
\item $PR \in [1.01,\,30]$ (when not fixed),
\item $T_{ref,1} \in [T_{c,in}-30,\,T_{c,in}-\Delta T_{pp,min,evap}]$,
\item $s_{ref,1}$ around saturated-vapor entropy at the cold-side inlet condition,
\item $T_{ref,3} \in [T_{h,in}+\Delta T_{pp,min,cond},\,T_{h,in}+150]$.
\end{itemize}

For each candidate, the model reconstructs all thermodynamic states, computes the required refrigerant mass flow from $\dot Q_{out,req}$, and rejects nonphysical or pinch-violating cycles. The best feasible solution maximizes $COP_{\mathrm{turb}}$ (equivalently minimizes $-COP_{\mathrm{turb}}$).

\subsection{Verification}

The implementation is verified against the model of \textbf{Martin T. White (PocketTHERM / SoftwareX)}. The reported agreement is within \textbf{0.01\% maximum relative error}. This supports correctness of the implemented cycle equations and property-coupled solution workflow.

\subsection{Cycle Input Case Used in Report (R1234ze(E))}

\begin{table}[h!]
\centering
\caption{\textbf{Cycle input parameters for the reported R1234ze(E) case}.}
\begin{tabular}{lll}
\hline
\textbf{Parameter} & \textbf{Value} & \textbf{Unit} \\
\hline
Refrigerant & R1234ze(E) & -- \\
$T_{h,in}$ & 353.15 & K \\
$T_{c,in}$ & 287.15 & K \\
$\dot m_h$ & 42 & kg/s \\
$c_{p,h}$ & 1885 & J/(kg\,K) \\
$\dot m_c$ & 40 & kg/s \\
$c_{p,c}$ & 1006 & J/(kg\,K) \\
$\eta_{\mathrm{compr}}$ & 0.78 & -- \\
$\eta_{\mathrm{turb}}$ & 0.87 & -- \\
$\eta_{\mathrm{shaft}}$ & 0.98 & -- \\
$\Delta T_{\mathrm{pp},1}$ & 10 & K \\
$\Delta T_{\mathrm{pp},3}$ & 10 & K \\
$\Delta T_{\mathrm{pp},4}$ & 10 & K \\
$\Delta T_{\mathrm{sh}}$ & 5 & K \\
\hline
\end{tabular}
\end{table}

\subsection{Cycle Input Case Used for CO2 Analyses}

\begin{table}[h!]
\centering
\caption{\textbf{Cycle input parameters used for CO2 analyses}.}
\begin{tabular}{lll}
\hline
\textbf{Parameter} & \textbf{Value} & \textbf{Unit} \\
\hline
Refrigerant & CO2 & -- \\
$T_{h,in}$ & 276 & K \\
$T_{c,in}$ & 240 & K \\
$\dot m_h$ & 42 & kg/s \\
$c_{p,h}$ & 1885 & J/(kg\,K) \\
$\dot m_c$ & 80 & kg/s \\
$c_{p,c}$ & 1006 & J/(kg\,K) \\
$\eta_{\mathrm{compr}}$ & 0.78 & -- \\
$\eta_{\mathrm{turb}}$ & 0.87 & -- \\
$\eta_{\mathrm{shaft}}$ & 0.98 & -- \\
$\Delta T_{\mathrm{pp},1}$ & 10 & K \\
$\Delta T_{\mathrm{pp},3}$ & 10 & K \\
$\Delta T_{\mathrm{pp},4}$ & 3 & K \\
$\Delta T_{\mathrm{sh}}$ & 5 & K \\
\hline
\end{tabular}
\end{table}

\subsection{Cycle COP Results (R1234ze(E) Case)}

\begin{table}[h!]
\centering
\caption{\textbf{Computed performance results for the R1234ze(E) case}.}
\begin{tabular}{lll}
\hline
\textbf{Metric} & \textbf{Value} & \textbf{Unit} \\
\hline
Optimized pressure ratio $PR$ & \textbf{12.684} & -- \\
$COP_{\mathrm{turb}}$ & \textbf{2.690} & -- \\
$COP_{\mathrm{is}}$ & 2.875 & -- \\
$COP_{\mathrm{isenth}}$ & 1.879 & -- \\
$\dot Q_{out}$ & 318492 & W \\
$\dot W_{compr}$ & 169529 & W \\
$\dot W_{turb}$ & 52171 & W \\
\hline
\end{tabular}
\end{table}

\textbf{Interpretation.} The model predicts a substantial COP uplift potential when turbine work recovery is included, while satisfying the imposed pinch constraints.
 in your main .tex document.

\section{Conceptual Heat Pump Model With Turbine Expansion}

\textbf{Objective.} The objective of this code is to quantify \textbf{how much system COP can increase} when replacing throttling losses with a \textbf{turbine expander} and mechanical work recovery, using a global optimization of cycle state variables.

\subsection{Modelling Approach}

The implementation solves a steady-state vapor-compression heat pump cycle with turbine expansion, pinch-point constraints, and performance post-processing. Thermophysical properties are evaluated with \textbf{CoolProp/REFPROP} through an AbstractState-based wrapper. The optimization is performed with \textbf{Differential Evolution} and candidate cycles are parameterized by \(PR,\,T_{ref,1},\,s_{ref,1},\,T_{ref,3}\) (or \(T_{ref,1},\,s_{ref,1},\,T_{ref,3}\) when \(PR\) is fixed).

\textbf{Key assumptions}
\begin{itemize}
\item Steady-state, 1D component model.
\item Fixed component efficiencies: $\eta_{\mathrm{compr}}$, $\eta_{\mathrm{turb}}$, and shaft efficiency $\eta_{\mathrm{shaft}}$.
\item Pinch constraints are enforced as \textbf{true minimum approach temperatures} along condenser and evaporator heat-transfer paths (not only at endpoints).
\item Refrigerant mass flow rate is determined from the required sink-side heat duty $\dot Q_{out,req}$ and candidate specific condenser heat rejection.
\item Subcritical and supercritical operating regions both handled, with a near-critical buffer to avoid numerically unreliable states.
\item Infeasible candidates receive a small positive penalty score, while feasible candidates are ranked by minimizing $-COP_{\mathrm{turb}}$.
\end{itemize}

\textbf{Block diagram (code workflow)}

\textit{(Requires in preamble: \texttt{\textbackslash usepackage\{tikz\}} and
\texttt{\textbackslash usetikzlibrary\{arrows.meta,positioning,shapes.geometric\}}.)}

\begin{figure}[h!]
\centering
\begin{tikzpicture}[
	node distance=1.25cm and 1.8cm,
	>=Latex,
	block/.style={draw, rounded corners, align=left, text width=0.28\textwidth, inner sep=6pt, font=\footnotesize},
	decision/.style={draw, diamond, aspect=2.0, align=center, text width=0.18\textwidth, inner sep=2pt, font=\footnotesize},
	line/.style={-Latex, thick}
]

\node[block] (in) {
\textbf{Inputs / configuration}\\
$\{T_{h,in},\,T_{c,in},\,\dot m_h,\,\dot m_c,\,c_{p,h},\,c_{p,c},\,\eta_{\mathrm{compr}},\,\eta_{\mathrm{turb}},\,\eta_{\mathrm{shaft}},\\
\Delta T_{\mathrm{pp},1},\,\Delta T_{\mathrm{pp},3},\,\Delta T_{\mathrm{pp},4},\,\Delta T_{\mathrm{sh}},\,\text{refrigerant}\}$
};

\node[decision, right=of in] (mode) {\textbf{Fixed $PR$?}};

\node[block, above right=0.35cm and 1.7cm of mode] (eval) {
\textbf{Cycle evaluation for candidate variables}\\
Solve stations $1\rightarrow2\rightarrow3\rightarrow4$:\\
$\{p_i,T_i,h_i,s_i,Q_i\}$ from CoolProp/REFPROP.
};

\node[block, below right=0.35cm and 1.7cm of mode] (obj) {
\textbf{Objective function $f(\mathbf{x})$}\\
$\mathbf{x}=[PR,\,T_{ref,1},\,s_{ref,1},\,T_{ref,3}]$ (or fixed-$PR$ subset),\\
$f(\mathbf{x})=-COP_{\mathrm{turb}}$ with infeasibility penalty.
};

\node[block, right=of eval] (checks) {
\textbf{Feasibility checks}\\
Pinch constraints: true condenser/evaporator $\Delta T_{\min}$ checks,\\
physical checks (quality/state validity),\\
near-critical buffer for robust optimization.
};

\node[block, right=of obj] (opt) {
\textbf{Global optimization (DE)}\\
Bounded search over candidate vector $\mathbf{x}$;\\
select $\mathbf{x}^*=\arg\max(COP_{\mathrm{turb}})$ among feasible solutions.
};

\node[block, right=of checks] (perf) {
\textbf{Performance post-processing}\\
Compute $\dot W_{\mathrm{compr}},\,\dot W_{\mathrm{turb}},\,\dot Q_{out},\\
COP_{\mathrm{turb}},\,COP_{\mathrm{is}},\,COP_{\mathrm{isenth}}$.
};

\node[block, below=1.1cm of perf] (out) {
\textbf{Outputs}\\
Best cycle state set, selected decision variables, COP metrics, and TS/PH plotting data.
};

\draw[line] (in) -- (mode);
\draw[line] (mode) -- node[above, font=\scriptsize] {yes} (eval);
\draw[line] (mode) -- node[below, font=\scriptsize] {no} (obj);
\draw[line] (eval) -- (checks);
\draw[line] (checks) -- (perf);
\draw[line] (obj) -- (opt);
\draw[line] (opt) |- (perf);
\draw[line] (perf) -- (out);

\end{tikzpicture}
\caption{\textbf{Detailed workflow of the turbine-assisted heat pump model and global multi-variable optimization.}}
\end{figure}

\subsection{Implemented Optimization Method (Current Version)}

The current solver uses \textbf{Differential Evolution} (global, population-based search) instead of local or pressure-ratio-only scanning. The candidate bounds are defined as:
\begin{itemize}
\item $PR \in [1.01,\,30]$ (when not fixed),
\item $T_{ref,1} \in [T_{c,in}-30,\,T_{c,in}-\Delta T_{pp,min,evap}]$,
\item $s_{ref,1}$ around saturated-vapor entropy at the cold-side inlet condition,
\item $T_{ref,3} \in [T_{h,in}+\Delta T_{pp,min,cond},\,T_{h,in}+150]$.
\end{itemize}

For each candidate, the model reconstructs all thermodynamic states, computes the required refrigerant mass flow from $\dot Q_{out,req}$, and rejects nonphysical or pinch-violating cycles. The best feasible solution maximizes $COP_{\mathrm{turb}}$ (equivalently minimizes $-COP_{\mathrm{turb}}$).

\subsection{Verification}

The implementation is verified against the model of \textbf{Martin T. White (PocketTHERM / SoftwareX)}. The reported agreement is within \textbf{0.01\% maximum relative error}. This supports correctness of the implemented cycle equations and property-coupled solution workflow.

\subsection{Cycle Input Case Used in Report (R1234ze(E))}

\begin{table}[h!]
\centering
\caption{\textbf{Cycle input parameters for the reported R1234ze(E) case}.}
\begin{tabular}{lll}
\hline
\textbf{Parameter} & \textbf{Value} & \textbf{Unit} \\
\hline
Refrigerant & R1234ze(E) & -- \\
$T_{h,in}$ & 353.15 & K \\
$T_{c,in}$ & 287.15 & K \\
$\dot m_h$ & 42 & kg/s \\
$c_{p,h}$ & 1885 & J/(kg\,K) \\
$\dot m_c$ & 40 & kg/s \\
$c_{p,c}$ & 1006 & J/(kg\,K) \\
$\eta_{\mathrm{compr}}$ & 0.78 & -- \\
$\eta_{\mathrm{turb}}$ & 0.87 & -- \\
$\eta_{\mathrm{shaft}}$ & 0.98 & -- \\
$\Delta T_{\mathrm{pp},1}$ & 10 & K \\
$\Delta T_{\mathrm{pp},3}$ & 10 & K \\
$\Delta T_{\mathrm{pp},4}$ & 10 & K \\
$\Delta T_{\mathrm{sh}}$ & 5 & K \\
\hline
\end{tabular}
\end{table}

\subsection{Cycle Input Case Used for CO2 Analyses}

\begin{table}[h!]
\centering
\caption{\textbf{Cycle input parameters used for CO2 analyses}.}
\begin{tabular}{lll}
\hline
\textbf{Parameter} & \textbf{Value} & \textbf{Unit} \\
\hline
Refrigerant & CO2 & -- \\
$T_{h,in}$ & 276 & K \\
$T_{c,in}$ & 240 & K \\
$\dot m_h$ & 42 & kg/s \\
$c_{p,h}$ & 1885 & J/(kg\,K) \\
$\dot m_c$ & 80 & kg/s \\
$c_{p,c}$ & 1006 & J/(kg\,K) \\
$\eta_{\mathrm{compr}}$ & 0.78 & -- \\
$\eta_{\mathrm{turb}}$ & 0.87 & -- \\
$\eta_{\mathrm{shaft}}$ & 0.98 & -- \\
$\Delta T_{\mathrm{pp},1}$ & 10 & K \\
$\Delta T_{\mathrm{pp},3}$ & 10 & K \\
$\Delta T_{\mathrm{pp},4}$ & 3 & K \\
$\Delta T_{\mathrm{sh}}$ & 5 & K \\
\hline
\end{tabular}
\end{table}

\subsection{Cycle COP Results (R1234ze(E) Case)}

\begin{table}[h!]
\centering
\caption{\textbf{Computed performance results for the R1234ze(E) case}.}
\begin{tabular}{lll}
\hline
\textbf{Metric} & \textbf{Value} & \textbf{Unit} \\
\hline
Optimized pressure ratio $PR$ & \textbf{12.684} & -- \\
$COP_{\mathrm{turb}}$ & \textbf{2.690} & -- \\
$COP_{\mathrm{is}}$ & 2.875 & -- \\
$COP_{\mathrm{isenth}}$ & 1.879 & -- \\
$\dot Q_{out}$ & 318492 & W \\
$\dot W_{compr}$ & 169529 & W \\
$\dot W_{turb}$ & 52171 & W \\
\hline
\end{tabular}
\end{table}

\textbf{Interpretation.} The model predicts a substantial COP uplift potential when turbine work recovery is included, while satisfying the imposed pinch constraints.
 in your main .tex document.

\section{Conceptual Heat Pump Model With Turbine Expansion}

\textbf{Objective.} The objective of this code is to quantify \textbf{how much system COP can increase} when replacing throttling losses with a \textbf{turbine expander} and mechanical work recovery, using a global optimization of cycle state variables.

\subsection{Modelling Approach}

The implementation solves a steady-state vapor-compression heat pump cycle with turbine expansion, pinch-point constraints, and performance post-processing. Thermophysical properties are evaluated with \textbf{CoolProp/REFPROP} through an AbstractState-based wrapper. The optimization is performed with \textbf{Differential Evolution} and candidate cycles are parameterized by \(PR,\,T_{ref,1},\,s_{ref,1},\,T_{ref,3}\) (or \(T_{ref,1},\,s_{ref,1},\,T_{ref,3}\) when \(PR\) is fixed).

\textbf{Key assumptions}
\begin{itemize}
\item Steady-state, 1D component model.
\item Fixed component efficiencies: $\eta_{\mathrm{compr}}$, $\eta_{\mathrm{turb}}$, and shaft efficiency $\eta_{\mathrm{shaft}}$.
\item Pinch constraints are enforced as \textbf{true minimum approach temperatures} along condenser and evaporator heat-transfer paths (not only at endpoints).
\item Refrigerant mass flow rate is determined from the required sink-side heat duty $\dot Q_{out,req}$ and candidate specific condenser heat rejection.
\item Subcritical and supercritical operating regions both handled, with a near-critical buffer to avoid numerically unreliable states.
\item Infeasible candidates receive a small positive penalty score, while feasible candidates are ranked by minimizing $-COP_{\mathrm{turb}}$.
\end{itemize}

\textbf{Block diagram (code workflow)}

\textit{(Requires in preamble: \texttt{\textbackslash usepackage\{tikz\}} and
\texttt{\textbackslash usetikzlibrary\{arrows.meta,positioning,shapes.geometric\}}.)}

\begin{figure}[h!]
\centering
\begin{tikzpicture}[
	node distance=1.25cm and 1.8cm,
	>=Latex,
	block/.style={draw, rounded corners, align=left, text width=0.28\textwidth, inner sep=6pt, font=\footnotesize},
	decision/.style={draw, diamond, aspect=2.0, align=center, text width=0.18\textwidth, inner sep=2pt, font=\footnotesize},
	line/.style={-Latex, thick}
]

\node[block] (in) {
\textbf{Inputs / configuration}\\
$\{T_{h,in},\,T_{c,in},\,\dot m_h,\,\dot m_c,\,c_{p,h},\,c_{p,c},\,\eta_{\mathrm{compr}},\,\eta_{\mathrm{turb}},\,\eta_{\mathrm{shaft}},\\
\Delta T_{\mathrm{pp},1},\,\Delta T_{\mathrm{pp},3},\,\Delta T_{\mathrm{pp},4},\,\Delta T_{\mathrm{sh}},\,\text{refrigerant}\}$
};

\node[decision, right=of in] (mode) {\textbf{Fixed $PR$?}};

\node[block, above right=0.35cm and 1.7cm of mode] (eval) {
\textbf{Cycle evaluation for candidate variables}\\
Solve stations $1\rightarrow2\rightarrow3\rightarrow4$:\\
$\{p_i,T_i,h_i,s_i,Q_i\}$ from CoolProp/REFPROP.
};

\node[block, below right=0.35cm and 1.7cm of mode] (obj) {
\textbf{Objective function $f(\mathbf{x})$}\\
$\mathbf{x}=[PR,\,T_{ref,1},\,s_{ref,1},\,T_{ref,3}]$ (or fixed-$PR$ subset),\\
$f(\mathbf{x})=-COP_{\mathrm{turb}}$ with infeasibility penalty.
};

\node[block, right=of eval] (checks) {
\textbf{Feasibility checks}\\
Pinch constraints: true condenser/evaporator $\Delta T_{\min}$ checks,\\
physical checks (quality/state validity),\\
near-critical buffer for robust optimization.
};

\node[block, right=of obj] (opt) {
\textbf{Global optimization (DE)}\\
Bounded search over candidate vector $\mathbf{x}$;\\
select $\mathbf{x}^*=\arg\max(COP_{\mathrm{turb}})$ among feasible solutions.
};

\node[block, right=of checks] (perf) {
\textbf{Performance post-processing}\\
Compute $\dot W_{\mathrm{compr}},\,\dot W_{\mathrm{turb}},\,\dot Q_{out},\\
COP_{\mathrm{turb}},\,COP_{\mathrm{is}},\,COP_{\mathrm{isenth}}$.
};

\node[block, below=1.1cm of perf] (out) {
\textbf{Outputs}\\
Best cycle state set, selected decision variables, COP metrics, and TS/PH plotting data.
};

\draw[line] (in) -- (mode);
\draw[line] (mode) -- node[above, font=\scriptsize] {yes} (eval);
\draw[line] (mode) -- node[below, font=\scriptsize] {no} (obj);
\draw[line] (eval) -- (checks);
\draw[line] (checks) -- (perf);
\draw[line] (obj) -- (opt);
\draw[line] (opt) |- (perf);
\draw[line] (perf) -- (out);

\end{tikzpicture}
\caption{\textbf{Detailed workflow of the turbine-assisted heat pump model and global multi-variable optimization.}}
\end{figure}

\subsection{Implemented Optimization Method (Current Version)}

The current solver uses \textbf{Differential Evolution} (global, population-based search) instead of local or pressure-ratio-only scanning. The candidate bounds are defined as:
\begin{itemize}
\item $PR \in [1.01,\,30]$ (when not fixed),
\item $T_{ref,1} \in [T_{c,in}-30,\,T_{c,in}-\Delta T_{pp,min,evap}]$,
\item $s_{ref,1}$ around saturated-vapor entropy at the cold-side inlet condition,
\item $T_{ref,3} \in [T_{h,in}+\Delta T_{pp,min,cond},\,T_{h,in}+150]$.
\end{itemize}

For each candidate, the model reconstructs all thermodynamic states, computes the required refrigerant mass flow from $\dot Q_{out,req}$, and rejects nonphysical or pinch-violating cycles. The best feasible solution maximizes $COP_{\mathrm{turb}}$ (equivalently minimizes $-COP_{\mathrm{turb}}$).

\subsection{Verification}

The implementation is verified against the model of \textbf{Martin T. White (PocketTHERM / SoftwareX)}. The reported agreement is within \textbf{0.01\% maximum relative error}. This supports correctness of the implemented cycle equations and property-coupled solution workflow.

\subsection{Cycle Input Case Used in Report (R1234ze(E))}

\begin{table}[h!]
\centering
\caption{\textbf{Cycle input parameters for the reported R1234ze(E) case}.}
\begin{tabular}{lll}
\hline
\textbf{Parameter} & \textbf{Value} & \textbf{Unit} \\
\hline
Refrigerant & R1234ze(E) & -- \\
$T_{h,in}$ & 353.15 & K \\
$T_{c,in}$ & 287.15 & K \\
$\dot m_h$ & 42 & kg/s \\
$c_{p,h}$ & 1885 & J/(kg\,K) \\
$\dot m_c$ & 40 & kg/s \\
$c_{p,c}$ & 1006 & J/(kg\,K) \\
$\eta_{\mathrm{compr}}$ & 0.78 & -- \\
$\eta_{\mathrm{turb}}$ & 0.87 & -- \\
$\eta_{\mathrm{shaft}}$ & 0.98 & -- \\
$\Delta T_{\mathrm{pp},1}$ & 10 & K \\
$\Delta T_{\mathrm{pp},3}$ & 10 & K \\
$\Delta T_{\mathrm{pp},4}$ & 10 & K \\
$\Delta T_{\mathrm{sh}}$ & 5 & K \\
\hline
\end{tabular}
\end{table}

\subsection{Cycle Input Case Used for CO2 Analyses}

\begin{table}[h!]
\centering
\caption{\textbf{Cycle input parameters used for CO2 analyses}.}
\begin{tabular}{lll}
\hline
\textbf{Parameter} & \textbf{Value} & \textbf{Unit} \\
\hline
Refrigerant & CO2 & -- \\
$T_{h,in}$ & 276 & K \\
$T_{c,in}$ & 240 & K \\
$\dot m_h$ & 42 & kg/s \\
$c_{p,h}$ & 1885 & J/(kg\,K) \\
$\dot m_c$ & 80 & kg/s \\
$c_{p,c}$ & 1006 & J/(kg\,K) \\
$\eta_{\mathrm{compr}}$ & 0.78 & -- \\
$\eta_{\mathrm{turb}}$ & 0.87 & -- \\
$\eta_{\mathrm{shaft}}$ & 0.98 & -- \\
$\Delta T_{\mathrm{pp},1}$ & 10 & K \\
$\Delta T_{\mathrm{pp},3}$ & 10 & K \\
$\Delta T_{\mathrm{pp},4}$ & 3 & K \\
$\Delta T_{\mathrm{sh}}$ & 5 & K \\
\hline
\end{tabular}
\end{table}

\subsection{Cycle COP Results (R1234ze(E) Case)}

\begin{table}[h!]
\centering
\caption{\textbf{Computed performance results for the R1234ze(E) case}.}
\begin{tabular}{lll}
\hline
\textbf{Metric} & \textbf{Value} & \textbf{Unit} \\
\hline
Optimized pressure ratio $PR$ & \textbf{12.684} & -- \\
$COP_{\mathrm{turb}}$ & \textbf{2.690} & -- \\
$COP_{\mathrm{is}}$ & 2.875 & -- \\
$COP_{\mathrm{isenth}}$ & 1.879 & -- \\
$\dot Q_{out}$ & 318492 & W \\
$\dot W_{compr}$ & 169529 & W \\
$\dot W_{turb}$ & 52171 & W \\
\hline
\end{tabular}
\end{table}

\textbf{Interpretation.} The model predicts a substantial COP uplift potential when turbine work recovery is included, while satisfying the imposed pinch constraints.
